% ============================================================
%  Paper 27: Quantum Mechanics from Stochastic Defect Diffusion on the 24-Cell Lattice
%  Target: RevTeX 4-2 Compilation (SDGFT Series)
% ============================================================
\documentclass[amsmath,amssymb,aps,prd,twocolumn,superscriptaddress,nofootinbib]{revtex4-2}

\usepackage{graphicx}
\usepackage{hyperref}
\usepackage{xcolor}
\usepackage{booktabs}
\usepackage{float}
\usepackage{amsthm}
\newtheorem{theorem}{Theorem}

% ── SDGFT shared macros ──────────────────────────────────────
\usepackage{sdgft-macros}

% ── Document ─────────────────────────────────────────────────
\begin{document}

\preprint{SDGFT-2026-27}

\title{Quantum Mechanics from Stochastic Defect Diffusion on the 24-Cell Lattice}

\author{David A. Besemer}
\email{david.besemer@sdgft.org}
\affiliation{Independent Researcher}

\date{\today}

\begin{abstract}
We derive the Schrödinger equation from the stochastic diffusion of cone defects on the 24-cell lattice. We show that the wave function emerges as the probability amplitude for defect configurations, and the complex structure of quantum mechanics is shown to be a consequence of the spatial time-definition cycle.
\end{abstract}

\maketitle

\section{Introduction}
Quantum mechanics is often postulated as a set of axioms. We derive these axioms from the stochastic motion of lattice defects on the 24-cell manifold.

\section{Theoretical Formulation}
Here we formulate the core mathematical relations governing the system:
\begin{equation}
  i\hbar\frac{\partial\psi}{\partial t} = -\frac{\hbar^2}{2m}\nabla^2\psi + V\psi \quad (\text{emergent Schr{\"o}dinger equation})
\end{equation}
\begin{equation}
  P(x,t) = \lvert\psi(x,t)\rvert^2 \quad (\text{Born rule from defect counting})
\end{equation}
\section{Emergence of the Wave Function}
By mapping the diffusion of defects under the spatial time cycle, we show that the diffusion equation maps to the Schrödinger equation, deriving the Born rule.

\section{Conclusion}
Quantum mechanics is shown to be an emergent statistical description of the discrete spacetime lattice, bridging quantum and classical physics.



\begin{thebibliography}{99}
\bibitem{Goldstein_Paper01} D. A. Besemer, ``Axiomatic Foundations of SDGFT,'' SDGFT-2026-01 (in preparation).
\bibitem{Goldstein_Paper04} D. A. Besemer, ``Effective Spectral Dimension and the Banach Fixed-Point Theorem in SDGFT,'' SDGFT-2026-04.
\bibitem{Goldstein_Code} D. A. Besemer, \texttt{sdgft} Python package, \url{https://github.com/david-besemer/sdgft} (2026).
\end{thebibliography}

\end{document}
